% Generated by scripts/assemble.py; corrections belong in source/corrections.json.
\examhead{141}{湖南大学 2024 年数学分析试卷}{134}
% p135-b22; source PDF page 135
\blockanchor{p135-b22}
\question{1}解答如下问题:\par
% p135-b23; source PDF page 135
\blockanchor{p135-b23}
(1) 已知  $0 < x_{1} < 1$ ,  $x_{n+1} = x_{n}(1 - x_{n})$ ,  $n = 1, 2, \cdots$ , 证明:  $\lim_{n \to \infty} nx_{n} = 1$ .\par
% p135-b24; source PDF page 135
\blockanchor{p135-b24}
(2) 设 $0 < p < 1, 0 < x_{1} < \frac{1}{p}, x_{n+1} = x_{n}(1 - px_{n})$, 证明: $\lim_{n \to \infty} nx_{n} = \frac{1}{p}$.\par
% p135-b25; source PDF page 135
\blockanchor{p135-b25}
\question{2}设函数 $f(x)$ 在点 $x = x_0$ 处可导.\par
% p135-b26; source PDF page 135
\blockanchor{p135-b26}
(1) 记 $x_{n} = f\left(x_{0} + \frac{1}{n^{2}}\right) + f\left(x_{0} + \frac{2}{n^{2}}\right) + \cdots + f\left(x_{0} + \frac{n}{n^{2}}\right) - nf(x_{0})$，证明：$\lim_{n \to \infty} x_{n} = \frac{1}{2} f'(x_{0})$.\par
% p136-b00; source PDF page 136
\blockanchor{p136-b00}
(2) 求极限  $\lim_{n\to\infty}\left(\sin\frac{1}{n^{2}}+\sin\frac{2}{n^{2}}+\cdots+\sin\frac{n}{n^{2}}\right)$ .\par
% p136-b01; source PDF page 136
\blockanchor{p136-b01}
(3) 求极限  $\lim_{n\to\infty}\left(1+\frac{1}{n^{2}}\right)\left(1+\frac{2}{n^{2}}\right)\cdots\left(1+\frac{n}{n^{2}}\right).$\par
% p136-b02; source PDF page 136
\blockanchor{p136-b02}
\question{3}解答如下问题:\par
% p136-b03; source PDF page 136
\par\Needspace{6\baselineskip}
\blockanchor{p136-b03}
(1) 设 $f(x)$ 为三次多项式, $x \in [-1, 1]$. 证明:\par
\nopagebreak[4]
% p136-b04; source PDF page 136
\blockanchor{p136-b04}
\displayformula{\int_ {- 1} ^ {1} f (x) \mathrm{d} x = \frac {1}{3} [ f (1) + 4 f (0) + f (- 1) ].}
% p136-b05; source PDF page 136
\par\Needspace{6\baselineskip}
\blockanchor{p136-b05}
(2) 设 $f(x)$ 为 $[a, b]$ 上的三次多项式, 证明:\par
\nopagebreak[4]
% p136-b06; source PDF page 136
\blockanchor{p136-b06}
\displayformula{\int_ {a} ^ {b} f (x) \mathrm{d} x = \frac {b - a}{6} \left[ f (a) + 4 f \left(\frac {a + b}{2}\right) + f (b) \right].}
% p136-b07; source PDF page 136
\blockanchor{p136-b07}
\question{4}设 $f(x)$ 在 $[0, +\infty)$ 上一致连续, 对任意的 $h > 0$, 序列 $\{f(nh)\}$ 极限存在. 证明: $\lim_{x \to +\infty} f(x)$ 存在.\par
% p136-b08; source PDF page 136
\blockanchor{p136-b08}
\question{5}设函数列 $\{f_n(x)\}$ 在 $[a, b]$ 上连续，且 $\{f_n(x)\}$ 在 $[a, b]$ 上一致收敛于 $f(x)$. 证明：\par
% p136-b09; source PDF page 136
\blockanchor{p136-b09}
(1) 存在 $M > 0$, 对任意的正整数 $n$ 及 $x \in [a, b]$, 有 $|f_n(x)| \leq M$, 且 $|f(x)| \leq M$.\par
% p136-b10; source PDF page 136
\blockanchor{p136-b10}
(2) 若 $g(x)$ 在 $(- \infty, + \infty)$ 上连续, 那么 $\{g(f_n(x))\}$ 在 $[a, b]$ 上一致收敛于 $g(f(x))$.\par
% p136-b11; source PDF page 136
\blockanchor{p136-b11}
\question{6}设 $f(x) = \frac{1}{4} x (2\pi - x), x \in [0, 2\pi]$.\par
% p136-b12; source PDF page 136
\blockanchor{p136-b12}
(1) 将 $f(x)$ 展开为 $[0, 2\pi]$ 上的 Fourier 级数, 并计算 $\sum_{n=1}^{\infty} \frac{1}{n^2}$.\par
% p136-b13; source PDF page 136
\blockanchor{p136-b13}
(2) 通过将 $f(x)$ 的 Fourier 级数逐项积分, 计算 $\sum_{n=1}^{\infty} \frac{1}{n^4}$.\par
% p136-b14; source PDF page 136
\blockanchor{p136-b14}
\question{7}设 $f(x), g(x)$ 在 $(a, b)$ 上连续, 什么情况下方程 $f(x)y = g(x)$ 在 $(a, b)$ 上确定了唯一的连续解?\par
% p136-b15; source PDF page 136
\blockanchor{p136-b15}
\question{8}设 $f(x)$ 在 $[0, +\infty)$ 上连续, 证明: 含参量积分 $\int_{0}^{+\infty} e^{-\alpha x} f(x) \mathrm{d}x$ 在 $\alpha \in [0, +\infty)$ 上一致收敛的充要条件为 $\int_{0}^{+\infty} f(x) \mathrm{d}x$ 收敛.\par
% p136-b16; source PDF page 136
\par\Needspace{6\baselineskip}
\blockanchor{p136-b16}
\question{9}计算曲面积分\par
\nopagebreak[4]
% p136-b17; source PDF page 136
\blockanchor{p136-b17}
\displayformula{\iint_ {S} x y z (x ^ {2} y ^ {2} + y ^ {2} z ^ {2} + z ^ {2} x ^ {2}) \mathrm{d} S.}
% p136-b18; source PDF page 136
\blockanchor{p136-b18}
其中 S 为  $x^{2} + y^{2} + z^{2} = a^{2} (a > 0)$  在第一卦限的部分.\par
